% Generated draft: requires mathematical and source review.
\documentclass[11pt]{article}
\usepackage[margin=25mm]{geometry}
\usepackage{fontspec}
\setmainfont{Latin Modern Roman}
\usepackage{amsmath,amssymb,mathtools,unicode-math}
\setmathfont{Latin Modern Math}
\usepackage{enumitem,needspace,xurl,hyperref}
\setlength{\parindent}{0pt}
\setlength{\parskip}{5pt}
\setlength{\emergencystretch}{3em}
\newcommand{\N}{\mathbb{N}}
\newcommand{\Z}{\mathbb{Z}}
\newcommand{\Q}{\mathbb{Q}}
\newcommand{\R}{\mathbb{R}}
\newcommand{\C}{\mathbb{C}}
\newcommand{\F}{\mathbb{F}}
\newcommand{\A}{\mathbb{A}}
\newcommand{\PP}{\mathbb{P}}
\newcommand{\E}{\mathbb{E}}
\newcommand{\eps}{\varepsilon}
\newcommand{\dd}{\,\mathrm{d}}
\begin{document}
\section*{UnsolvedMath ID: 10600006}
\textbf{Source problem code:} catalogue entry
\subsection{Definitions and mathematical statement}
A welded knot is a virtual-knot diagram modulo the classical and virtual Reidemeister moves and one chosen forbidden move, here the upper forbidden move. Allowing both forbidden moves is reported in the source to trivialize virtual knots. The source says the upper move preserves the Wirtinger fundamental group and quandle; it also says every welded knot is the closure of a welded braid. Exact local pictures and orientation conventions are not included.

The first request is to construct effective distinguishing invariants of welded knots and, if possible, classify them. No completeness or computational criterion for ``good'' invariants is specified.

Let \(\operatorname{Tube}\) denote Satoh's map from welded knots or links to ambient-isotopy classes of ribbon torus embeddings in four-dimensional space. The source asks whether this map is faithful, that is, whether
\[
 \operatorname{Tube}(K_1)\cong\operatorname{Tube}(K_2)
 \quad\Longrightarrow\quad K_1\sim_{\rm welded}K_2
\]
for the intended knot/link classes. The torus orientation, coorientation, ambient space and equivalence conventions are absent; they matter to the precise injectivity question. The source reports preservation of the combinatorial fundamental group by an isomorphism with the complement group of the torus embedding. Its parenthetical ``in fact the quandle'' must be checked as a corresponding quandle statement, not read as identifying a quandle with a group.

Finally, the source asks whether a welded knot with ``trivial fundamental group'' is a trivial welded knot. Literally, if \(G(K)\) is the usual Wirtinger group, the wording asks
\[
 G(K)=\{1\}\ \Longrightarrow\ K\sim_{\rm welded}U.
\]
But the usual one-component Wirtinger group has abelianization \(\Z\), and the unknot group is \(\Z\), not the one-element group. The intended condition may therefore be ``group isomorphic to the unknot group.'' That would instead ask whether \(G(K)\cong\Z\) implies \(K\sim_{\rm welded}U\). These formulations are distinguished rather than silently interchanged.

The supplied partial-progress summary reports invariants, restricted injectivity results and kernel work for the Tube map, while calling general faithfulness and the trivial-group characterization unresolved. It does not identify the restricted classes or settle the equivalence conventions, and its bibliographic identifiers contain placeholders. Its claims remain supplied assessments requiring primary-source verification.
\subsection{Short English statement}
Develop and classify welded-knot invariants, ask whether Satoh's ribbon-torus map detects welded equivalence, and clarify whether an unknot-like fundamental group forces welded triviality. The record's literal trivial-group wording and unspecified Tube-map conventions are substantial formulation issues. The supplied progress summary has incomplete references.
\subsection{Sources}
\begin{itemize}
\item Local UnsolvedMath registry, permanent ID \(10600006\), exact source code \texttt{catalogue entry}, ``Virtual-knot problem 6 --- Welded Knots.'' No individual proposer or proposed year is supplied.
\item \emph{Unsolved Problems in Virtual Knot Theory and Combinatorial Knot Theory}, revised survey, dedicated problem-list item 6: \url{https://arxiv.org/abs/1409.2823}. Background: 60 top-level items, compound items retained intact; source-TeX extraction, accessed 2026-07-29. Several in-statement reference slots are empty.
\item Supplied research references, with incomplete identifiers retained literally: Audoux, ``On the welded Tube map'' (2014), arXiv:1408.xxxx; ``Some generalizations of Satoh's Tube map'' (2021), arXiv:2103.xxxx; ``Non-triviality of welded knots and ribbon torus-knots'' (2024), arXiv:2403.xxxx; ``An unknotting invariant for welded knots'' (2020), arXiv:2008.xxxx. These identifiers are placeholders, not complete retrievable citations.
\item Further supplied incomplete references: homotopy classification of ribbon tubes and welded string links (2014), arXiv:1407.xxxx; ``Bridge numbers and meridional ranks of knotted surfaces and welded knots'' (2021), arXiv:2111.xxxx; ``Welded graphs, Wirtinger groups and knotted punctured spheres'' (2023), arXiv:2311.xxxx. The summary ends in a truncated ``related'' remark; no missing text is supplied.
\item Registry status partially\_solved and classification PARTIAL-PROGRESS report progress on all parts, restricted Tube-map injectivity, and no definitive general faithfulness/trivial-group resolution. Source evidence NEEDS\_REVIEW presents an open problem and requires a current primary-literature audit.
\item Dataset attribution: Ulam AI / UnsolvedMath Contributors, CC BY 4.0. Source rights NEEDS\_REVIEW: publicly accessible author/publisher copy, redistribution terms requiring review.
\end{itemize}
\end{document}
